%% kapa: exam=1 round=23 year=2012 period=1 form=A subject=civil_law subject_ko=민법 seq=15 number=15
%% source: Q-Net 감정평가사 기출문제 (cst00309) artlSeq=5052578 file=raw/1008/5052578/2012년 제23회 감정평가사 제1차시험 기출문제(1교시 A형).pdf
%% conversion: kordoc 4.24.51 -> tools/gen1.py (auto); review=proofread-sonnet; flags=none
\begin{question}{15}
甲은 乙에게 500만 원을 빌리면서 800만 원 상당의 명품시계에 질권을 설정하였다. 그 후 乙은 丙으로부터 300만 원을 빌리면서 甲의 시계에 관하여 丙에게 다시 질권을 설정하여 주고 그 사실을 甲에게 통지하였다. 다음 설명 중 甲이 질권의 소멸을 주장하여 시계를 돌려받을 수 있는 경우는?
\begin{choices}
  \item 甲이 乙에게 500만 원을 변제한 경우
  \item 甲이 丙에게 300만 원을 변제한 경우
  \item 甲이 丙에게 200만 원을 변제한 경우
  \item 甲이 丙에게 300만 원을 변제하고 乙에게 200만 원을 변제한 경우
  \item 甲이 丙에게 200만 원을 변제하고 乙에게 300만 원을 변제한 경우
\end{choices}
\end{question}
\answer{4}
